
\subsubsection{Bin Configurations}
\label{subsec:bp-binconfig}

The configuration-based formulation provides a common
starting point for the approximation techniques developed
by both de la Vega--Lueker and Karmarkar--Karp.

A \textbf{bin configuration} describes the composition
of a single bin, specifying the number of items of
each distinct size packed into it. It can also be
viewed as a multiset of item sizes whose total size
does not exceed the unit capacity. Two configurations
are considered identical if they contain the same
number of items of each distinct size.

In the general Bin Packing problem, the number of
possible configurations can grow exponentially with
the number of items. This is because each configuration
corresponds to a non-empty sub-multiset of the input
items, although not every sub-multiset is necessarily
feasible. Therefore, the total number of feasible
configurations is at most $2^n-1$, where $n$ is the
number of input items.


\subsubsection{Integer Linear Programming Formulation}
\label{subsubsec:bp-ilp}

Consider a Bin Packing instance $I$ containing $n$
items of $m$ distinct sizes $y_1,\ldots,y_m$, with
$n_j$ items of size $y_j$. Let $\mathcal{C}=\{C_1,\ldots,C_q\}$ denote the
set of all feasible non-empty bin configurations. For each configuration $C_i$, let $c_{ij}$
denote the number of items of size $y_j$
contained in that configuration. By feasibility,
$\sum_{j=1}^{m} c_{ij}y_j \leq 1,\, \forall i\in\{1,\ldots,q\}.$ 
We represent a packing by a vector
$\mathbf{a}=(a_1,\ldots,a_q)$, where $a_i$
denotes the number of bins assigned to
configuration $C_i$. The Bin Packing problem can then be formulated as the following Integer Linear Program (ILP):

\begin{equation}
\label{eq:bp-ilp}
\begin{aligned}
    \min_{\mathbf{a}} \quad
        & \sum_{i=1}^{q} a_i \\[4pt]
    \text{subject to} \quad
        & \sum_{i=1}^{q} a_i c_{ij} \geq n_j,
          && \forall j\in\{1,\ldots,m\},\\
        & a_i \in \mathbb{Z}_{\geq 0},
          && \forall i\in\{1,\ldots,q\}.
\end{aligned}
\end{equation}

The objective minimizes the total number of
bins used. For each item size $y_j$, the
corresponding constraint ensures that the
selected configurations provide sufficient
capacity to accommodate all $n_j$ items of
that size.

Each feasible integer vector
$\mathbf{a}=(a_1,\ldots,a_q)$ corresponds to
a feasible packing of $I$, where $a_i$ bins
are assigned to configuration $C_i$.
Any excess item slots can be left unused.

For a feasible solution produced by an algorithm,
we denote the number of bins used as
$
    \ALG(I)=\sum_{i=1}^{q}a_i.
$
Thus, minimizing the ILP objective is equivalent
to minimizing the number of bins required to pack
$I$, and its optimal value is $\OPT(I)$.

However, as discussed previously in 
\Cref{subsec:bp-binconfig}, the number of
feasible configurations $q$ can be exponential
in $n$. Consequently, even constructing the
complete ILP may require exponentially many
variables. This motivates considering a restricted version
of Bin Packing in which the number of feasible
configurations can be bounded more tightly.


\subsubsection{Restricted Bin Packing: Bounded Configurations}
\label{subsubsec:bounded-config}

Consider a restricted Bin Packing instance in which
every item has size at least $\varepsilon>0$.
Since each bin has unit capacity, it can contain
at most $k=\lfloor 1/\varepsilon\rfloor$ items. Under this restriction, the number of feasible
configurations can be bounded in terms of $k$
and the number of distinct item sizes $m$.

\begin{lemma}
\label{lem:bounded-config}
For any Bin Packing instance with $m$ distinct
item sizes, each at least $\varepsilon>0$,
the number of non-empty feasible bin configurations
$q$ satisfies
$
    q\leq \binom{k+m}{m}-1,
    \, k=\lfloor 1/\varepsilon\rfloor.
$
\end{lemma}

\begin{proof}
Let $c_j$ denote the number of items of size
$y_j$ in a configuration. Since each bin contains
at most $k$ items,
$
    \sum_{j=1}^{m}c_j\leq k,
    \, c_j\in\mathbb{Z}_{\geq 0}.
$

Introducing a slack variable $c_R\geq 0$ gives
$
    \sum_{j=1}^{m}c_j+c_R=k.
$

By the stars-and-bars formula, this equation
has $\binom{k+m}{m}$ non-negative integer
solutions. Excluding the empty configuration
and configurations violating the capacity
constraint yields the desired bound.
\end{proof}

A detailed combinatorial illustration of this
counting argument is provided in
\Cref{app:bounded-config}. Consequently, for fixed $k$ and $m$, the configuration ILP in \Cref{subsubsec:bp-ilp}
contains a bounded number of variables,
independent of the total number of items $n$.


\subsubsection{Linear Programming Relaxation}

Although restricting the number of configurations
reduces the size of the ILP formulation, solving
the resulting integer program optimally remains
computationally challenging. Hence, A natural approach is to relax the integrality constraints $a_i \in \mathbb{Z}_{\geq 0}$ to $a_i \geq 0$, obtaining a Linear Programming
(LP) relaxation.

Both de la Vega--Lueker and Karmarkar--Karp
exploit configuration-based formulations and
LP relaxations, but they employ rounding techniques to
construct feasible integral packings from the relaxed 
LP solution. In the following sections, we examine their
respective techniques and analyze their
approximation guarantees and running-time
complexities.


\paragraph{De La Vega--Lueker's Approach}

Let $\mathbf{a}^*=(a_1^*,\ldots,a_q^*)$ be an optimal
basic feasible solution to the LP relaxation, with
objective value $LIN(I)=\sum_{i=1}^{q}a_i^*$. 
Since the LP has $m$ covering constraints, any
basic feasible solution has at most $m$ non-zero
variables. This property allows us to construct
an integral solution by rounding up each variable,
while introducing at most $m$ additional bins.

\textbf{Rounding up the LP solution.}
We construct an integer vector
$\hat{\mathbf{a}}=(\hat{a}_1,\ldots,\hat{a}_q)$
by setting
\[
    \hat{a}_i=\lceil a_i^*\rceil,
    \qquad \forall i\in\{1,\ldots,q\}.
\]

Since $\hat{a}_i\geq a_i^*$ for every $i$,
the rounded solution continues to satisfy all
covering constraints:
\[
    \sum_{i=1}^{q}\hat{a}_i c_{ij}
    \geq
    \sum_{i=1}^{q}a_i^*c_{ij}
    \geq n_j,
    \qquad \forall j\in\{1,\ldots,m\}.
\]

Hence, $\hat{\mathbf{a}}$ is a feasible integer
solution to the ILP. Moreover, each non-zero
variable increases by less than one during
rounding. Since there are at most $m$ such
variables, the total number of bins increases
by at most $m$.


\textbf{Greedy Deflation.}
Although the rounded solution satisfies the
covering constraints, it may contain more items
of a particular size than required. We therefore perform a polynomial-time \textbf{deflation} procedure that removes excess items by replacing selected configurations with smaller feasible
configurations. A detailed batch-deflation algorithm, together
with an illustrative example and correctness
argument, is provided in \Cref{app:greedy-deflation}.

By rounding up each of the $q$ LP variables,
the number of bins increases by at most $q$.
Since deflation does not introduce additional
bins, we obtain $\ALG(I) \leq \LIN(I)+q \leq\OPT(I)+q.$

\paragraph{Karmarkar--Karp's approach}

Karmarkar and Karp proposed a slightly different ways of rounding optimal solution $\mathbf{a}^*$ of the LP relaxation. Then, a following packed is constructed by seperate each of components $a_i^*$ of $\mathbf{a}^*$ into two parts -- the integral part $\lfloor a_i^* \rfloor$ and the fractional part $f_i^*$ where $a_i^* = f_i^* + \lfloor a_i^* \rfloor$. Let $\mathbf{p}^* = \{\lfloor a_1^* \rfloor, \dots , \lfloor a_q^* \rfloor \}$, and $\mathbf{f}^* = \{f_1^*, \dots , f_n^* \}$, obviously, $\mathbf{a}^* = \mathbf{p}^* + \mathbf{f}^*$ and $\mathbf{p}^*$ can represent some valid packing of some parts of the current instance $I$. Karmarkar and Karp remove this packing represented by the integral parts from the original packing where $n_j' = n_j - \sum_{i = 1}^{q}\lfloor a_i^* \rfloor c_{ij}, \, \forall j \in \{1, \dots, m\}$. 

The remaining items will form a new instance $I'$ which satisfies following inequalites $n_j' \leq \sum_{i = 1}^{q} f_i^*$, and we can infer that $\mathbf{f}^*$ is an optimal solution to the LP relaxation $I'$ by optimal substructure. To pack the remaining items, there are two possible ways of feasible packing: (1) a packing with exactly $m$ bins is the feasible solution, (2) we follow a greedy algorithm and in the previous part, we have proved that $\ALG(I') = 2\SIZE(I') + 1$.

To prove the first way of packing is feasible, we started from the foundamental property of the LP where an optimal solution can be attained at a vertex of the feasible region, provided such a vertex exists. In our formulation, the non-negativity constraints and the existence of a finite optimum guarantee that an optimal vertex solution exists. Karmarkar and Karp tries to make $\mathbf{a}^*$ an optimal vertex solution. Since we have $m$ constraints (excluding the nonnegativity requirements) and if we expressed the constraints in the standard form of matrix multiplication $A\mathbf{x} \geq \mathbf{b}$, $A$ has $m$ rows. Given the \Cref{lem:lp_nonzero_variables}, since the number of nonzero variables is at most $m$, one bin of config $i$ for each nonzero $f_i^*$ is a valid packing. The second way of proving is demostrated in the greedy proof which has been presented before.

Obviously, both packing will have a running time polynomial of the input instance. Karmarkar and Karp will try both packing and return the minimum one as final packing $\ALG(I') \leq \min(m, 2\SIZE(I') + 1) \leq \frac{m + 2\SIZE(I') + 1}{2} \leq \SIZE(I') + \frac{m + 1}{2}$. The total packing will thus be expressed as $\ALG(I) = \sum_i{\mathbf{p}^*} + \ALG(I').$ In addition, for any instance $I$, we can have $\LIN(I) \geq \SIZE(I)$ because $\LIN(I) = \sum_{i=1}^{q}{a_i} \geq \sum_{i = 1}^{q}{a_i}\sum_{j = 1}^{m}a_ic_{ij} = \sum_{j = 1}^{m}y_j\sum_{i = 1}^q{a_i c_{ij}} \geq \sum_{i = 1}^m{y_j n_j} = \SIZE(I)$.

we can hence derive following relationship: $\OPT(I) \leq \ALG(I) = \sum_i{\mathbf{p}^*} + \ALG(I') \leq \LIN(I) - \LIN(I') + \ALG(I') \leq \LIN(I) - \LIN(I') + \SIZE(I') + \frac{m + 1}{2} \leq \LIN(I) - \LIN(I') + \LIN(I') + \frac{m + 1}{2} = \LIN(I) + \frac{m + 1}{2}.$